\documentclass[12pt]{article}
\def\activityid{F07-U-v2}\usepackage[letterpaper,margin=.65in,top=.60in,bottom=.65in]{geometry}
\usepackage[T1]{fontenc}\usepackage[default]{sourcesanspro}
\usepackage{microtype,tikz,array,tabularx,fancyhdr,amsmath}\usepackage[hidelinks]{hyperref}
\usetikzlibrary{arrows.meta,calc}
\setlength{\parindent}{0pt}\setlength{\parskip}{7pt}\setlength{\headheight}{0pt}\setlength{\footskip}{26pt}\setlength{\emergencystretch}{2em}
\pagestyle{fancy}\fancyhf{}\renewcommand{\headrulewidth}{0pt}\renewcommand{\footrulewidth}{.4pt}
\fancyfoot[L]{\scriptsize Bellingham Math Circle / Week 7 / \activityid}\fancyfoot[R]{\scriptsize \thepage}
\definecolor{ink}{gray}{.12}\definecolor{soft}{gray}{.95}\color{ink}
\newcommand{\PageTitle}[2]{{\fontsize{10}{12}\selectfont\bfseries WEEK 7\quad GAME LAB\hfill #1\par}\vspace{3pt}{\fontsize{26}{29}\selectfont\bfseries #2\par}\vspace{2pt}}
\newcommand{\NameLine}{{\small Name: \rule{2.65in}{.4pt}\hfill Date: \rule{1.35in}{.4pt}\par}\vspace{4pt}}
\newcommand{\Task}[2]{\par\vspace{3pt}{\large\bfseries #1}\par #2\par}
\newcommand{\Subhead}[1]{\par\vspace{3pt}{\large\bfseries #1}\par}
\newcommand{\RuleBox}[1]{\par\begingroup\setlength{\fboxsep}{8pt}\colorbox{soft}{\parbox{\dimexpr\linewidth-16pt\relax}{#1}}\endgroup\par}
\newcommand{\WriteBox}[2]{\par\textbf{#1}\par\vspace{2pt}\begin{tikzpicture}\draw[black!40,line width=.5pt] (0,0) rectangle (\dimexpr\linewidth-.5pt\relax,#2);\end{tikzpicture}\par}
\newcommand{\StudentFooter}{\vfill{\small Build first. Explain by pointing, talking, or drawing.}\par}
\newcommand{\Code}[1]{\texttt{#1}}

\hypersetup{pdftitle={Week 7: grades-4-5}}
\begin{document}
\PageTitle{GRADES 4--5}{A pattern is a conjecture}\NameLine
\RuleBox{Two players alternate. Take \textbf{1, 3, or 4} counters each turn. Taking 2 is illegal. Never take more than remain. Taking the last counter wins.}\Task{1. Build the evidence.}{Play from 10 and 14. Then label positions W or L from the viewpoint of the player about to move. Put L at zero.}
\begin{center}\renewcommand{\arraystretch}{1.85}\begin{tabular}{p{.68in}|p{.68in}|p{.68in}|p{.68in}|p{.68in}|p{.68in}|p{.68in}}0&1&2&3&4&5&6\\\hline&&&&&&\\7&8&9&10&11&12&13\\\hline&&&&&&\\14&15&16&17&18&19&20\\\hline&&&&&&\\21&22&23&24&25&26&27\\\hline&&&&&&\\\end{tabular}\end{center}
\Task{2. Make every label earn its place.}{For W, give one legal move to L. For L, check all legal moves go to W. Find a disputed position and defend it.}
\WriteBox{What repeat do you conjecture? State it exactly.}{1in}
\textbf{Next question:} Is seeing four repeats a proof that the pattern continues forever? What feature of the rules could force it?\StudentFooter
\newpage
\PageTitle{GRADES 4--5}{Prove the repeat forever}\NameLine
\RuleBox{Two players alternate. Take \textbf{1, 3, or 4} counters each turn. Taking 2 is illegal. Never take more than remain. Taking the last counter wins.}\Task{3. Two obligations.}{Your proposed losing piles must satisfy \textbf{both} tests: (a) no legal move goes from one proposed L to another; (b) every other positive pile has a legal move to a proposed L.}
\Task{4. Use groups of seven.}{Write a pile as some whole groups of seven plus a remainder from 0 to 6. Complete the table using your conjecture.}\begin{center}\renewcommand{\arraystretch}{2}\begin{tabular}{c|c|p{1.25in}|p{1.6in}}Remainder&W or L&Take this many&New remainder\\\hline0&&&\\1&&&\\2&&&\\3&&&\\4&&&\\5&&&\\6&&&\end{tabular}\end{center}
\WriteBox{Why do the two obligations prove a winning strategy?}{.9in}
\textbf{Check the boundary:} Your suggested move must be legal even for the smallest positive pile in that remainder class. Why does a decreasing pile eventually reach zero?\StudentFooter
\newpage
\PageTitle{GRADES 4--5}{Change a rule; challenge a proof}\NameLine
\RuleBox{New game: take \textbf{1, 3, or 5}. Two players alternate. Taking the last counter wins. This is not ``take any number from 1 to 5.''}
\Task{5. Predict before playing.}{Who can force a win from 6? From 7? From 100? Test small piles, then find one explanation for every pile.}
\WriteBox{My new losing piles and both parts of my proof:}{1.5in}
\Task{6. How much memory does the old game need?}{Return to moves 1, 3, 4. Copy a stretch of your labels onto a straight strip. Put a paper frame around \textbf{four consecutive W/L labels}. Which three cells decide the next label? Add that label, then slide the frame one cell right. Find a matching window elsewhere and show why both next labels must agree.}
\WriteBox{Why identical four-label windows have identical futures:}{.95in}
\textbf{Go further:} A four-label window has at most \(2^4=16\) patterns. Explain why some window must recur. Does this prove repetition starts at zero, or only that it eventually starts?
\StudentFooter
\end{document}
